\documentclass[aspectratio=169,10pt]{beamer}
\usetheme[numbering=fraction,progressbar=frametitle]{metropolis}
\usepackage{fontspec}
\setsansfont{Latin Modern Sans}
\usepackage{amsmath,amssymb,mathtools,graphicx,booktabs}
\usepackage{polyglossia}
\setdefaultlanguage{french}
\setbeamertemplate{navigation symbols}{}
\setbeamerfont{frametitle}{size=\large}
\setbeamerfont{normal text}{size=\normalsize}
\setlength{\parskip}{.45em}
\hypersetup{pdfauthor={},pdfsubject={Lecture slides},pdfkeywords={},pdftitle={Divisibilité et congruences}}
\providecommand{\N}{\mathbb{N}}
\providecommand{\Z}{\mathbb{Z}}
\providecommand{\D}{\mathbb{D}}
\providecommand{\Q}{\mathbb{Q}}
\providecommand{\R}{\mathbb{R}}
\providecommand{\C}{\mathbb{C}}
\providecommand{\vect}[1]{\overrightarrow{#1}}
\providecommand{\vi}{\vec{\imath}}
\providecommand{\vj}{\vec{\jmath}}
\providecommand{\vk}{\vec{k}}
\providecommand{\norm}[1]{\left\lVert #1\right\rVert}
\providecommand{\vc}[2]{\begin{pmatrix}#1\\#2\end{pmatrix}}
\providecommand{\deter}[2]{\det\left(#1,#2\right)}
\providecommand{\abs}[1]{\left\lvert #1\right\rvert}
\providecommand{\dd}{\,\mathrm{d}}
\providecommand{\e}{\mathrm{e}}
\providecommand{\Card}{\operatorname{Card}}
\providecommand{\E}{\operatorname{E}}
\providecommand{\V}{\operatorname{V}}
\providecommand{\degre}{^{\circ}}
\providecommand{\un}[1]{\,\mathrm{#1}}
\providecommand{\sci}[2]{#1\times10^{#2}}
\providecommand{\pgcd}{\operatorname{PGCD}}
\providecommand{\ppcm}{\operatorname{PPCM}}
\providecommand{\llbracket}{\mathopen{[\mkern-3mu[}}
\providecommand{\rrbracket}{\mathclose{]\mkern-3mu]}}
\title{Divisibilité et congruences}
\subtitle{Cours et exemples}
\date{}
\author{}
\begin{document}
\begin{frame}[plain]
\vfill
{\usebeamerfont{title}\inserttitle\par}
\medskip
{\usebeamerfont{subtitle}\insertsubtitle\par}
\medskip
{\usebeamercolor[fg]{progress bar}\rule{\textwidth}{.5pt}}
\vfill
\end{frame}
\begin{frame}{Objectifs}
Utiliser la divisibilité dans \(\Z\) et ses propriétés ; effectuer une division euclidienne ; calculer avec des congruences ; établir un critère de divisibilité ; étudier des restes de puissances.
\end{frame}
\begin{frame}{Activité d'introduction}
Le 17 octobre 2026 est un samedi. On numérote les jours de la semaine de \(0\) (samedi) à \(6\) (vendredi) : dimanche porte le numéro \(1\), lundi le numéro \(2\), etc.
\end{frame}
\begin{frame}{Activité d'introduction (2)}
\begin{enumerate}
\setcounter{enumi}{0}\item Quel jour de la semaine sera-t-il \(7\) jours plus tard ? \(14\) jours ? \(700\) jours ?

\item Quel jour sera-t-il \(100\) jours plus tard ? Expliquer comment obtenir la réponse sans compter les jours un à un.

\item Même question pour \(1\,000\) jours.

\item L'année 2027 n'est pas bissextile. Quel jour de la semaine sera le 17 octobre 2027 ?

\item Le 1er janvier 2027 arrive \(76\) jours après le 17 octobre 2026. Quel jour sera-t-il ?

\end{enumerate}
\end{frame}
\begin{frame}{Activité d'introduction (3)}
\begin{enumerate}
\setcounter{enumi}{5}\item Deux durées \(a\) et \(b\), en jours, mènent au même jour de la semaine. Que peut-on dire de \(a-b\) ? Formuler une condition générale.

\end{enumerate}
\end{frame}
\begin{frame}{Divisibilité dans \(\Z\)}
Soit \(a\) et \(b\) deux entiers relatifs. On dit que \(a\) \emph{divise} \(b\), et on note \(a\mid b\), s'il existe un entier relatif \(k\) tel que \(b=ka\). On dit aussi que \(a\) est un \emph{diviseur} de \(b\), ou que \(b\) est un \emph{multiple} de \(a\).
\end{frame}
\begin{frame}{Divisibilité dans \(\Z\)}
Soit \(a\), \(b\), \(c\) des entiers relatifs.

\begin{enumerate}
\setcounter{enumi}{0}
\item Si \(a\mid b\) et \(b\mid c\), alors \(a\mid c\).

\item Si \(a\mid b\) et \(a\mid c\), alors \(a\mid bu+cv\) pour tous entiers relatifs \(u\) et \(v\).

\item Si \(a\mid b\) et \(b\neq0\), alors \(\abs a\le\abs b\).

\end{enumerate}
\end{frame}
\begin{frame}{Divisibilité dans \(\Z\)}
1. \(b=ka\) et \(c=k'b\) donnent \(c=(k'k)a\). 2. \(b=ka\) et \(c=k'a\) donnent \(bu+cv=(ku+k'v)a\). 3. \(b=ka\) avec \(k\neq0\), donc \(\abs k\ge1\) et \(\abs b=\abs k\abs a\ge\abs a\).
\end{frame}
\begin{frame}{Divisibilité dans \(\Z\)}
Soit \(a\in\Z\) et \(b\in\N\), \(b\ge1\). Il existe un unique couple d'entiers \((q;r)\) tel que \(a=bq+r\) et \(0\le r<b\). On appelle \(q\) le \emph{quotient} et \(r\) le \emph{reste}.
\end{frame}
\begin{frame}{Divisibilité dans \(\Z\)}
\emph{Existence.} Posons \(q=\left\lfloor\dfrac ab\right\rfloor\), le plus grand entier inférieur ou égal à \(\dfrac ab\). Alors \(q\le\dfrac ab<q+1\), donc \(bq\le a<bq+b\), et \(r=a-bq\) vérifie \(0\le r<b\).

\emph{Unicité.} Si \(a=bq+r=bq'+r'\) avec \(0\le r,r'<b\), alors \(b(q-q')=r'-r\) et \(-b<r'-r<b\). Le seul multiple de \(b\) strictement compris entre \(-b\) et \(b\) est \(0\), donc \(r=r'\) puis \(q=q'\).
\end{frame}
\begin{frame}{Divisibilité dans \(\Z\)}
Le reste est toujours positif : \(-47=5\times(-10)+3\), donc la division de \(-47\) par \(5\) a pour quotient \(-10\) et pour reste \(3\), et non \(-9\) et \(-2\).
\end{frame}
\begin{frame}{Congruences}
Soit \(n\ge2\) un entier. Deux entiers relatifs \(a\) et \(b\) sont congrus modulo \(n\), ce que l'on note \(a\equiv b\pmod n\), si \(n\mid a-b\).
\end{frame}
\begin{frame}{Congruences}
\(a\equiv b\pmod n\) si et seulement si \(a\) et \(b\) ont le même reste dans la division euclidienne par \(n\). En particulier, tout entier est congru modulo \(n\) à son reste, qui appartient à \(\llbracket 0;n-1\rrbracket\).
\end{frame}
\begin{frame}{Congruences}
Écrivons \(a=nq+r\) et \(b=nq'+r'\) avec \(0\le r,r'<n\). Alors \(a-b=n(q-q')+(r-r')\) et \(-n<r-r'<n\). Donc \(n\mid a-b\) si et seulement si \(n\mid r-r'\), si et seulement si \(r=r'\).
\end{frame}
\begin{frame}{Congruences}
Si \(a\equiv b\pmod n\) et \(c\equiv d\pmod n\), alors
\[a+c\equiv b+d,\qquad a-c\equiv b-d,\qquad ac\equiv bd,\qquad a^k\equiv b^k\pmod n\]
pour tout entier naturel \(k\). La relation est aussi transitive : si \(a\equiv b\) et \(b\equiv c\pmod n\), alors \(a\equiv c\pmod n\).
\end{frame}
\begin{frame}{Congruences}
Somme : \((a+c)-(b+d)=(a-b)+(c-d)\) est somme de deux multiples de \(n\). Produit : \(ac-bd=a(c-d)+d(a-b)\), multiple de \(n\). Puissance : récurrence sur \(k\) à partir du produit. Transitivité : \(a-c=(a-b)+(b-c)\).
\end{frame}
\begin{frame}{Congruences}
On ne simplifie pas une congruence comme une égalité : \(2\times3\equiv2\times1\pmod4\), mais \(3\not\equiv1\pmod4\). La division n'est possible que dans des cas précis (chapitre 7).
\end{frame}
\begin{frame}{Congruences}
On cherche une puissance congrue à \(1\) (ou à \(-1\)), puis on utilise la division euclidienne de l'exposant. Exemple : \(2^3=8\equiv1\pmod7\) et \(2\,026=3\times675+1\), donc \(2^{2026}=\left(2^3\right)^{675}\times2\equiv1^{675}\times2=2\pmod7\).
\end{frame}
\begin{frame}{Congruences}
Pour étudier une expression modulo \(n\), on dresse un tableau donnant sa valeur modulo \(n\) pour chaque reste possible de \(x\). Exemple : modulo \(4\), \(x\equiv0,1,2,3\) donne \(x^2\equiv0,1,0,1\). Un carré est donc congru à \(0\) ou \(1\) modulo \(4\).
\end{frame}
\begin{frame}{Congruences}
Si \(N=\sum_{k=0}^{p}a_k10^k\) s'écrit avec les chiffres \(a_0,\dots,a_p\), comme \(10\equiv1\pmod9\), on a \(10^k\equiv1\) et \(N\equiv a_0+a_1+\dots+a_p\pmod9\). Un entier est divisible par \(9\) si et seulement si la somme de ses chiffres l'est.
\end{frame}
\end{document}
